% Part: incompleteness
% Chapter: introduction
% Section: decidability

\documentclass[../../../include/open-logic-section]{subfiles}

\begin{document}

\olfileid{inc}{int}{dec}

\olsection{অণির্ণেয়তা ও অসম্পূর্ণতা}

গ্যোডেলের অসম্পূর্ণতা উপপাদ্যের প্রমাণে syntax-এর arithmetization দরকার। কিন্তু
তা ছাড়াও, কোনো তত্ত্ব সব !!{decidable} সম্বন্ধকে !!{represents} করে—শুধু এই
অনুমান থেকে কিছু সুন্দর ফল পাওয়া যায়। প্রমাণটি halting problem-এর
অণির্ণেয়তার প্রমাণের মতোই একটি diagonal argument।

\begin{thm}
যদি $\Gamma$ একটি সঙ্গতিপূর্ণ তত্ত্ব হয় এবং প্রতিটি !!{decidable} সম্বন্ধকে
!!{represents} করে, তবে $\Gamma$ !!{decidable} নয়।
\end{thm}

\begin{proof}
ধরি $\Gamma$ !!{decidable}। আমরা দেখাব, $\Gamma$ যদি প্রত্যেক
!!{decidable} সম্বন্ধকে !!{represents} করে, তবে এটি অসঙ্গতিপূর্ণ হতে বাধ্য।

!!^{decidable} ধর্মগুলি (এক-পদী সম্বন্ধগুলি) একটি মুক্ত চলরাশি-যুক্ত
!!{formula}s দিয়ে প্রকাশ করা যায়। $!A_0(x)$, $!A_1(x)$, \dots-কে এই ধরনের
সব !!{formula}s-এর একটি গণনসাধ্য তালিকা ধরি। এখন নিচের সেট $D \subseteq \Nat$
বিবেচনা করি:
\[
D = \Setabs{n}{\Gamma \Proves \lnot !A_n(\num{n})}
\]
$D$ !!{decidable}, কারণ $n \in D$ কি না পরীক্ষা করতে প্রথমে $!A_n(x)$
গণনা করা যায়, এবং সেখান থেকে $\lnot !A_n(\num{n})$ গঠন করা যায়।
$\num{n}$ পদটি প্রতিটি মুক্ত $x$-এর ঘটনার স্থানে $!A_n(x)$-এ বসানো এবং
$!A_n(\num{n})$-এর সামনে $\lnot$ বসানো যান্ত্রিক কাজ।
$\Gamma$ !!{decidable} হওয়ার অনুমান থেকে পরীক্ষা করা যায়
$\lnot !A_n(\num{n}) \in \Gamma$ কি না। থাকলে $n \in D$, না থাকলে
$n \notin D$। তাই $D$-ও !!{decidable}।
% BN-SRC-203 TARGET-CORRECTION-BEGIN
\par\smallskip\noindent\textnormal{\small\textbf{উৎস-সংশোধন (BN-SRC-203):} হিমায়িত প্রমাণে diagonal পরিবারের $!A_n$-এর সাবস্ক্রিপ্ট পরের দুটি ব্যবহারে হারিয়ে গিয়ে $!A(\num{n})$ লেখা হয়েছে। বাংলা পাঠে উভয় স্থানে $!A_n(\num{n})$ রাখা হয়েছে, যাতে পরীক্ষিত সূত্রটি একই পরিবারের $n$-তম সদস্য থাকে।} % BN-SRC-203 TARGET-CORRECTION-END

যেহেতু $\Gamma$ সব !!{decidable} ধর্মকে !!{represents} করে, তাই এটি
!!{represents}~$D$। $D$-কে !!{represents}s করা !!{formula}s-গুলি $\Gamma$-র
মধ্যেই $!A_0(x)$, $!A_1(x)$, \dots তালিকায় আছে। অতএব এমন একটি সংখ্যা $d$ নিই
যে $!A_d(x)$ !!{represents} $D$-কে $\Gamma$-তে। যদি $d \notin D$ হয়,
তবে $!A_d(x)$ $D$-কে !!{represents} করার কারণে
$\Gamma \Proves \lnot !A_d(\num{d})$। কিন্তু এতে $d$-এর $D$-র সংজ্ঞায়িত
শর্ত পূরণ হয়, তাই $d \in D$। এটি $d \notin D$-র বিরোধী। অতএব পরোক্ষ
প্রমাণে $d \in D$।

$d \in D$ হওয়ায় $D$-র সংজ্ঞা থেকে
$\Gamma \Proves \lnot !A_d(\num{d})$। অন্য দিকে, $!A_d(x)$ $D$-কে
!!{represents} করে $\Gamma$-তে বলে
$\Gamma \Proves !A_d(\num{d})$। অতএব $\Gamma$ অসঙ্গতিপূর্ণ।
\end{proof}

\begin{explain}
আগের উপপাদ্য দেখায়, সব !!{decidable} সম্বন্ধকে !!{represents}s করা কোনো
সঙ্গতিপূর্ণ তত্ত্ব !!{decidable} হতে পারে না। আমরা দেখাব, $\Th{Q}$ সব
!!{decidable} সম্বন্ধকে !!{represents} করে; ফলে $\Th{Q}$-কে অন্তর্ভুক্ত করা
সব তত্ত্ব—যেমন $\Th{PA}$ ও $\Th{TA}$—ও তা করে এবং তাই !!{decidable} নয়।
(এই তত্ত্বগুলি standard model-এ সত্য হওয়ায় সবগুলিই সঙ্গতিপূর্ণ।)

এই ফল ব্যবহার করে প্রথম অসম্পূর্ণতা উপপাদ্যের একটি দুর্বল রূপও পাওয়া যায়।
যে কোনো !!{axiomatizable} ও !!{complete} তত্ত্ব !!{decidable}। সুতরাং যে
সঙ্গতিপূর্ণ !!{axiomatizable} তত্ত্ব সব !!{decidable} ধর্মকে !!{represents}s
করে, তা
!!{complete} হতে পারে না।
\end{explain}

\begin{thm}
যদি $\Gamma$ !!{axiomatizable} ও !!{complete} হয়, তবে এটি !!{decidable}।
\end{thm}

\begin{proof}
যে কোনো অসঙ্গতিপূর্ণ তত্ত্ব !!{decidable}, কারণ অসঙ্গতিপূর্ণ তত্ত্বে সব
!!{sentence}s থাকে; তাই ``$!A \in \Gamma$ কি?'' প্রশ্নের উত্তর সব সময়
``হ্যাঁ'', অর্থাৎ তা সিদ্ধান্ত করা যায়।

তাই ধরি $\Gamma$ সঙ্গতিপূর্ণ এবং একই সঙ্গে !!{axiomatizable} ও !!{complete}।
$\Gamma$ !!{axiomatizable} হওয়ায় এটি !!{computably enumerable}। কারণ
$\Gamma$-র স্বতঃসিদ্ধ থেকে সব সঠিক !!{derivation}s একটি গণনসাধ্য অপেক্ষক দিয়ে
তালিকাভুক্ত করা যায়। সঠিক !!{derivation} থেকে সেটি যে !!{sentence}
!!{derive}s তা গণনা করা যায়; অতএব একটি গণনসাধ্য অপেক্ষক $\Gamma$-র সব
উপপাদ্য তালিকাভুক্ত করে। একটি
!!{sentence} তখনই $\Gamma$-এর উপপাদ্য যখন $\lnot !A$ উপপাদ্য নয়; কারণ $\Gamma$ সঙ্গতিপূর্ণ
ও !!{complete}। তাই $!A \in \Gamma$
কি না এভাবে সিদ্ধান্ত করা যায়: $\Gamma$-র উপপাদ্যগুলি তালিকাভুক্ত করতে থাকি।
তালিকায় $!A$ এলে জানি $\Gamma \Proves !A$; আর $\lnot !A$ এলে জানি
$\Gamma \Proves/ !A$। $\Gamma$ !!{complete} বলে এই দুটির একটি অবশেষে আসবেই,
তাই পদ্ধতিটি শেষ পর্যন্ত উত্তর দেয়।
\end{proof}

\begin{cor}
\ollabel{cor:incompleteness}
যদি $\Gamma$ সঙ্গতিপূর্ণ, !!{axiomatizable}, এবং প্রতিটি !!{decidable} ধর্মকে
!!{represents} করে, তবে এটি !!{complete} নয়।
\end{cor}

\begin{proof}
$\Gamma$ !!{complete} হলে আগের উপপাদ্য অনুসারে (কারণ এটি !!{axiomatizable}
ও সঙ্গতিপূর্ণ) !!{decidable} হতো। কিন্তু $\Gamma$ সব !!{decidable} ধর্মকে
!!{represents} করায় প্রথম উপপাদ্য অনুসারে এটি !!{decidable} নয়।
\end{proof}

\begin{prob}
দেখাও যে $\Th{TA} = \Setabs{!A}{\Sat{N}{!A}}$ !!{axiomatizable} নয়। ধরে
নিতে পারো যে $\Th{TA}$ সব decidable ধর্মকে প্রতিনিধিত্ব করে।
\end{prob}

যেমন, $\Th{Q}$ সব !!{decidable} ধর্মকে !!{represents}s করে—এটি প্রতিষ্ঠা করতে
পারলে, উপপাদ্যটি বলে $\Th{Q}$ অসম্পূর্ণ। কিন্তু এই প্রমাণ কোনো স্বাধীন
!!{sentence}-এর উদাহরণ দেয় না; কেবল দেখায় যে এমন !!a{sentence} থাকবেই।
এর জন্য syntax arithmetize করে গ্যোডেলের মূল প্রমাণের ধারণা অনুসরণ করতে হবে।
এবং প্রথম দাবিটিও এখনও দেখাতে হবে—$\Th{Q}$ সত্যিই সব !!{decidable} ধর্মকে
!!{represents} করে।

মনে রাখা দরকার, প্রত্যেক \emph{আগ্রহজনক} তত্ত্ব অসম্পূর্ণ বা অণির্ণেয় নয়।
কিছু তত্ত্ব গ্যোডেলের ফল প্রয়োগের শর্ত পূরণ না করেও আকর্ষণীয় গাণিতিক তথ্য
প্রকাশ করতে যথেষ্ট শক্তিশালী। যেমন,
$\Th{Pres} = \Setabs{!A \in \Lang{L_{A^+}}}{\Sat{N}{!A}}$ হলো
standard model-এ সত্য, $\times$-চিহ্নহীন পাটিগণিতের ভাষার !!{sentence}s-গুলির সেট;
এটি একই সঙ্গে পূর্ণ ও decidable। এই তত্ত্বকে Presburger arithmetic বলা
হয় এবং এটি কেবল $\Obj{0}$, $\prime$ ও $+$ দিয়ে প্রকাশযোগ্য প্রাকৃতিক সংখ্যা
সম্বন্ধে সব সত্য প্রমাণ করে।

\end{document}
